% !TeX program = pdflatex
% PROJET ML & IA - DOCUMENT SCIENTIFIQUE EVOLUTIF
% Importer le dossier ZIP dans Overleaf, choisir main.tex comme fichier principal.
\documentclass[11pt,a4paper,oneside]{report}
% ================================================================
% CONFIGURATION GLOBALE : éditer principalement les lignes ci-dessous
% Compatible Overleaf gratuit / compilateur pdfLaTeX / sans options payantes
% ================================================================
\usepackage[T1]{fontenc}
\usepackage[utf8]{inputenc}
\usepackage[french]{babel}
\usepackage{lmodern}
\usepackage{microtype}
\usepackage[a4paper,left=27mm,right=25mm,top=27mm,bottom=26mm,headheight=16pt]{geometry}
\usepackage{graphicx}
\usepackage{tikz}
\usetikzlibrary{calc,positioning,arrows.meta}
\usepackage[most]{tcolorbox}
\usepackage{amsmath,amssymb}
\usepackage{booktabs,tabularx,array,longtable}
\usepackage{enumitem}
\usepackage{titlesec}
\usepackage{fancyhdr}
\usepackage{lastpage}
\usepackage{caption}
\usepackage{float}
\usepackage{hyperref}
\usepackage{url}
\usepackage{xcolor}

% --- Identité du document : MODIFIER CES CHAMPS ---
\newcommand{\Ecole}{École Centrale Casablanca}
\newcommand{\Cours}{Machine Learning et Intelligence Artificielle}
\newcommand{\TitreRapport}{Autoencodeurs pour le clustering\\et la détection d'anomalies}
\newcommand{\SousTitreRapport}{Application à la surveillance vibratoire des roulements}
\newcommand{\MembreA}{Prénom NOM du membre A}
\newcommand{\MembreB}{Prénom NOM du membre B}
\newcommand{\MembreC}{Prénom NOM du membre C}
\newcommand{\MembreD}{Prénom NOM du membre D}
\newcommand{\Encadrant}{À compléter}
\newcommand{\AnneeProjet}{2026}
\newcommand{\PeriodeProjet}{22 octobre -- 30 novembre 2026}
\newcommand{\VersionRapport}{0.1}
\newcommand{\DateMiseAJour}{27 octobre 2026}

% --- Logos et photo FACULTATIFS ---
% Ajouter ces fichiers dans le dossier images/, sans toucher au code.
\newcommand{\CheminLogoECC}{images/logo_ecc.png}
\newcommand{\CheminLogoPartenaire}{images/logo_partenaire.png}
\newcommand{\CheminPhotoBanc}{images/banc_essai.jpg}

% --- Palette sobre, académique, industrielle ---
\definecolor{bleunuit}{HTML}{14253D}
\definecolor{bleufonce}{HTML}{223C59}
\definecolor{turquoise}{HTML}{0E8794}
\definecolor{orclair}{HTML}{DBAD54}
\definecolor{grisclair}{HTML}{F3F6F8}
\definecolor{grisfonce}{HTML}{596A77}
\definecolor{bleupale}{HTML}{EAF3F5}
\definecolor{orangealerte}{HTML}{A36B21}

\hypersetup{pdftitle={Autoencodeurs pour le clustering et la détection d'anomalies},
            pdfauthor={Équipe de projet, École Centrale Casablanca},
            colorlinks=true,linkcolor=bleufonce,citecolor=turquoise,urlcolor=turquoise}
\setlength{\parindent}{0pt}
\setlength{\parskip}{0.55em}
\renewcommand{\arraystretch}{1.18}
\setlist[itemize]{leftmargin=*,topsep=3pt,itemsep=2pt}
\setlist[enumerate]{leftmargin=*,topsep=3pt,itemsep=2pt}
\captionsetup{font=small,labelfont={bf,color=bleunuit}}

% --- Titres et navigation ---
\titleformat{\chapter}[display]
  {\sffamily\bfseries\color{bleunuit}}
  {\textcolor{turquoise}{\small\MakeUppercase{Chapitre}\ \thechapter}}
  {3pt}{\Huge}
\titlespacing*{\chapter}{0pt}{-6pt}{20pt}
\titleformat{\section}{\Large\sffamily\bfseries\color{bleufonce}}{\thesection}{0.8em}{}
\titleformat{\subsection}{\normalsize\sffamily\bfseries\color{turquoise}}{\thesubsection}{0.8em}{}

\pagestyle{fancy}
\fancyhf{}
\fancyhead[L]{\small\sffamily\color{grisfonce}Projet Autoencodeurs $\cdot$ ML \& IA}
\fancyhead[R]{\small\sffamily\color{turquoise}Document évolutif $\vert$ v\VersionRapport}
\fancyfoot[L]{\footnotesize\color{grisfonce}\Ecole}
\fancyfoot[R]{\footnotesize\color{grisfonce}\thepage}
\renewcommand{\headrulewidth}{0.3pt}
\renewcommand{\footrulewidth}{0pt}
\fancypagestyle{plain}{\fancyhf{}\fancyfoot[L]{\footnotesize\color{grisfonce}\Ecole}\fancyfoot[R]{\footnotesize\color{grisfonce}\thepage}\renewcommand{\headrulewidth}{0pt}}

% --- Blocs d'information réutilisables ---
\newtcolorbox{encadreinfo}[1][]{enhanced,breakable,colback=bleupale,
  colframe=turquoise!45,boxrule=0.35pt,arc=1.5mm,
  borderline west={2.5pt}{0pt}{turquoise},left=4mm,right=4mm,top=3mm,bottom=3mm,#1}
\newtcolorbox{encadreattente}[1][]{enhanced,breakable,colback=grisclair,
  colframe=grisfonce!20,boxrule=0.35pt,arc=1mm,
  borderline west={2pt}{0pt}{orclair},left=4mm,right=4mm,top=3mm,bottom=3mm,#1}
\newcommand{\acompleter}[1]{\begin{encadreattente}\textbf{À compléter au cours du projet.}\\ #1\end{encadreattente}}


\begin{document}
\pagenumbering{roman}
% ================================================================
% PAGE DE GARDE : automatiquement fonctionnelle sans logos ni photos
% ================================================================
\begin{titlepage}
\newgeometry{left=21mm,right=21mm,top=17mm,bottom=17mm}
\thispagestyle{empty}
\begin{tikzpicture}[remember picture,overlay]
  \fill[bleunuit] (current page.north west) rectangle
    ([yshift=-9.35cm]current page.north east);
  \fill[turquoise] ([yshift=-9.35cm]current page.north west) rectangle
    ([yshift=-9.24cm]current page.north east);
  \fill[orclair] ([xshift=0cm,yshift=-9.35cm]current page.north west) rectangle
    ([xshift=3.9cm,yshift=-9.24cm]current page.north west);
  \draw[white,opacity=0.10,line width=17pt]
    ($(current page.north east)+(-0.9cm,-1.45cm)$) circle (2.25cm);
  \draw[turquoise,opacity=0.25,line width=1pt]
    ($(current page.north east)+(-2.5cm,-6.2cm)$) circle (2.5cm);
  \draw[white,opacity=0.19,line width=1.1pt]
    ($(current page.north west)+(0.0cm,-7.0cm)$) --
    ++(2.0cm,0cm) -- ++(0.35cm,-0.38cm) -- ++(0.5cm,0.77cm)
    -- ++(0.7cm,-1.07cm) -- ++(0.6cm,0.64cm) -- ++(1cm,0cm)
    -- ++(0.45cm,-0.3cm) -- ++(0.65cm,0.30cm) -- ++(2cm,0cm);
\end{tikzpicture}

\noindent
\begin{minipage}[t]{.73\linewidth}
  \IfFileExists{\CheminLogoECC}{\includegraphics[height=13mm,keepaspectratio]{\CheminLogoECC}}
  {\fcolorbox{white}{white}{\parbox[c][12mm][c]{39mm}{\centering\sffamily\bfseries\color{bleunuit}\small ÉCOLE CENTRALE\\CASABLANCA}}}
\end{minipage}%
\begin{minipage}[t]{.27\linewidth}
  \raggedleft
  \IfFileExists{\CheminLogoPartenaire}{\includegraphics[height=13mm,keepaspectratio]{\CheminLogoPartenaire}}
  {\fcolorbox{white}{white}{\parbox[c][12mm][c]{30mm}{\centering\sffamily\color{grisfonce}\scriptsize Logo\\facultatif}}}
\end{minipage}

\vspace{8mm}
{\sffamily\bfseries\color{orclair}\fontsize{10}{13}\selectfont
 PROJET ACADÉMIQUE $\cdot$ MACHINE LEARNING \& IA\par}
\vspace{3mm}
{\color{white}\sffamily\bfseries\fontsize{25}{30}\selectfont
 \TitreRapport\par}
\vspace{3mm}
{\color{white!86}\sffamily\fontsize{12}{17}\selectfont
 \SousTitreRapport\par}
\vspace{18mm}
{\color{turquoise}\footnotesize\sffamily\bfseries
 DOSSIER SCIENTIFIQUE ÉVOLUTIF \quad $\vert$ \quad VERSION \VersionRapport\par}

\vspace{15mm}
\begin{minipage}[t]{.57\linewidth}
  {\Large\sffamily\bfseries\color{bleunuit}Équipe projet}\par\vspace{2mm}
  \begin{tabular}{@{}>{\sffamily\bfseries\color{turquoise}}l@{\hspace{6mm}}l@{}}
    A & \MembreA\\[1mm]
    B & \MembreB\\[1mm]
    C & \MembreC\\[1mm]
    D & \MembreD
  \end{tabular}
\end{minipage}%
\hfill
\begin{minipage}[t]{.38\linewidth}
  {\Large\sffamily\bfseries\color{bleunuit}Encadrement}\par\vspace{2mm}
  {\color{grisfonce}\small\Encadrant}\par\vspace{7mm}
  {\color{turquoise}\sffamily\bfseries\small COURS}\par
  {\small\Cours}
\end{minipage}

\vspace{11mm}
\begin{tcolorbox}[colback=grisclair,colframe=grisclair,boxrule=0pt,
     arc=1mm,left=5mm,right=5mm,top=3mm,bottom=3mm]
\begin{tabularx}{\linewidth}{@{}>{\bfseries\color{bleunuit}}p{48mm}X@{}}
  Période du projet & \PeriodeProjet\\
  Responsable rapport & Membre B, avec contributions de toute l'équipe\\
  Dernière mise à jour & \DateMiseAJour
\end{tabularx}
\end{tcolorbox}
\vfill
{\color{grisfonce}\small\sffamily\Ecole\hfill\AnneeProjet}
\end{titlepage}
\restoregeometry


\chapter*{Note aux lecteurs}
\addcontentsline{toc}{chapter}{Note aux lecteurs}
\begin{encadreinfo}
\textbf{Nature du document :} version de cadrage scientifique, \emph{non définitive}.\\
\textbf{Statut :} premières sections rédigées, démarche expérimentale proposée, résultats non disponibles.\\
\textbf{Mise à jour :} \DateMiseAJour{} $\vert$ version \VersionRapport.
\end{encadreinfo}
Ce document accompagne l'avancement du projet et sert de support de travail commun aux quatre membres de l'équipe. Il sera enrichi et corrigé à mesure que la description du banc d'essai, les métadonnées d'acquisition, les traitements et les résultats seront vérifiés.

Les affirmations sur les \textbf{données disponibles} sont distinguées des \textbf{choix méthodologiques envisagés}. Aucune performance de modèle ni aucun résultat expérimental ne sont annoncés avant leur mesure. Les sections encore incomplètes sont conservées volontairement afin de guider la rédaction progressive.

\tableofcontents
\clearpage
\pagenumbering{arabic}

\chapter{Introduction générale et cadrage}
\label{chap:introduction}

\section{Contexte industriel}
La disponibilité et la fiabilité des équipements constituent des enjeux importants pour la continuité de la production, la sécurité et les coûts d'exploitation. Sur les machines tournantes, une dégradation mécanique peut perturber le fonctionnement avant d'aboutir à une défaillance. Les démarches de surveillance conditionnelle cherchent à exploiter des mesures représentatives de l'état de l'équipement pour orienter les opérations de maintenance.

Parmi ces mesures, les vibrations offrent une information privilégiée sur le comportement dynamique des organes mécaniques. Les changements de rigidité, les jeux, les chocs et les défauts de surface peuvent modifier les caractéristiques des signaux vibratoires. L'analyse des vibrations ne se réduit cependant pas à un seul indicateur : elle mobilise des mesures, des traitements du signal et une connaissance du fonctionnement de la machine \cite{boulenger2009}.

Les roulements occupent une place centrale dans de nombreuses machines tournantes. Une anomalie touchant un élément roulant ou une bague peut entraîner des phénomènes impulsionnels et modifier la répartition temporelle et fréquentielle des vibrations. L'exploitation de ces signatures dépend néanmoins du régime de fonctionnement, des conditions d'acquisition et de la transmission des vibrations jusqu'aux capteurs.

\section{Apport potentiel de l'apprentissage non supervisé}
L'analyse vibratoire conventionnelle repose notamment sur des indicateurs statistiques et des représentations spectrales. Elle fournit des informations utiles à l'interprétation mécanique, mais le choix des caractéristiques pertinentes peut être délicat lorsque les signaux sont complexes, bruités ou acquis dans plusieurs conditions de fonctionnement.

Les \textbf{autoencodeurs} constituent une famille de réseaux neuronaux entraînés à reconstruire leurs données d'entrée à partir d'une représentation intermédiaire. Leur usage est étudié, entre autres, pour apprendre des caractéristiques de machines tournantes \cite{shao2017} et pour détecter des anomalies à partir de l'erreur de reconstruction de séries temporelles multicapteurs \cite{malhotra2016}.

Dans ce projet, deux finalités distinctes seront examinées :
\begin{itemize}
  \item le \textbf{clustering}, pour explorer les regroupements naturels des observations dans un espace de caractéristiques appris ;
  \item la \textbf{détection d'anomalies}, pour étudier la possibilité de distinguer des états défectueux à partir d'un modèle de référence associé au fonctionnement sain.
\end{itemize}

Ces finalités ne doivent pas être confondues avec une classification supervisée : les étiquettes de défauts pourront servir à évaluer les regroupements, mais ne devront pas être utilisées pour entraîner la représentation dans un protocole déclaré non supervisé.

\section{Problématique scientifique}
Les données disponibles couvrent plusieurs états de roulements, plusieurs régimes de fonctionnement et deux voies de mesure d'accélération verticale. Une modification du signal peut donc résulter d'un défaut, d'un changement de régime ou de conditions expérimentales. Cette diversité constitue à la fois un intérêt scientifique et une source possible de confusion pour les algorithmes.

\begin{encadreinfo}
\textbf{Problématique proposée :}
\emph{Dans quelle mesure les autoencodeurs permettent-ils d'apprendre, à partir de mesures vibratoires expérimentales, des représentations utiles au regroupement des états de roulements et à la détection d'anomalies, tout en maîtrisant l'influence des conditions de fonctionnement ?}
\end{encadreinfo}

Cette problématique sera réexaminée après l'analyse exploratoire et la clarification des métadonnées d'acquisition.

\section{Objectifs du projet}
\textbf{Objectif général.} Étudier et évaluer des approches fondées sur les autoencodeurs pour le clustering et la détection d'anomalies dans des signaux vibratoires de roulements.

\textbf{Objectifs spécifiques envisagés :}
\begin{enumerate}
  \item caractériser la structure, la qualité et la variabilité des mesures expérimentales ;
  \item identifier des représentations pertinentes des signaux et concevoir des autoencodeurs adaptés ;
  \item analyser la structure des représentations latentes au moyen de méthodes de clustering ;
  \item mesurer la capacité d'un modèle entraîné sur un état sain à signaler des acquisitions anormales ;
  \item comparer ces solutions à des approches de référence plus simples ;
  \item évaluer la robustesse des résultats aux différents régimes et voies de mesure ;
  \item préparer une démonstration reproductible de la solution retenue.
\end{enumerate}

\section{Périmètre et prudence terminologique}
Les acquisitions proviennent d'un \textbf{banc d'essai expérimental de roulements}. Le dispositif ne doit pas être présenté comme un suivi direct d'une flotte ferroviaire en exploitation sans preuve documentaire.

Le terme \emph{maintenance prédictive} est employé ici pour situer la finalité industrielle générale. Le projet porte d'abord sur la \textbf{surveillance d'état et le diagnostic de comportements anormaux}. Il ne vise pas, en l'état actuel des informations, à estimer une durée de vie résiduelle (RUL), car l'existence de trajectoires de dégradation temporelle n'est pas établie.

\chapter{Cadre scientifique et état de l'art préliminaire}
\label{chap:cadre}

\section{Principes de l'analyse vibratoire}
Une vibration mécanique peut être caractérisée par son amplitude, sa fréquence et la nature des excitations qui la produisent. Les grandeurs cinématiques usuelles sont le déplacement, la vitesse et l'accélération. Pour des phénomènes rapides et impulsionnels, la mesure d'accélération peut être particulièrement informative, sous réserve que la chaîne de mesure possède une bande passante adaptée \cite{boulenger2009}.

Le manuel de Boulenger et Pachaud aborde notamment les grandeurs vibratoires (chapitre 1), le traitement des signaux (chapitre 2), les stratégies de surveillance (chapitre 4) et les défauts de roulements (section 5.13). Ces éléments constituent des \emph{pistes à examiner}, et non une liste d'outils à employer obligatoirement.

\section{Indicateurs temporels classiques}
Pour une fenêtre numérique de $N$ observations $x_1,\ldots,x_N$, on peut définir la valeur efficace
\begin{equation}
  x_{\mathrm{RMS}}=\sqrt{\frac{1}{N}\sum_{i=1}^{N}x_i^2}.
  \label{eq:rms}
\end{equation}
Une définition courante du facteur de crête est
\begin{equation}
  C_f=\frac{\max_i |x_i|}{x_{\mathrm{RMS}}},
  \label{eq:crest}
\end{equation}
et le kurtosis de Pearson peut s'écrire
\begin{equation}
  K=\frac{\frac{1}{N}\sum_{i=1}^{N}(x_i-\bar{x})^4}
          {\left[\frac{1}{N}\sum_{i=1}^{N}(x_i-\bar{x})^2\right]^2},
  \quad \bar{x}=\frac{1}{N}\sum_{i=1}^{N}x_i.
  \label{eq:kurtosis}
\end{equation}
Ces indicateurs ne suffisent pas, pris isolément, à identifier une pathologie mécanique. Ils pourront toutefois être étudiés comme \textbf{caractéristiques explicatives}, \textbf{points de comparaison} ou \textbf{variables d'exploration}.

\section{Analyse fréquentielle : condition préalable}
Si la fréquence d'échantillonnage $f_s$ est connue et l'acquisition adéquate, la transformée de Fourier discrète permet d'étudier les contributions fréquentielles du signal. Pour une fenêtre de $N$ échantillons, l'espacement nominal des fréquences discrètes est
\begin{equation}
  \Delta f=\frac{f_s}{N}.
  \label{eq:resolution}
\end{equation}
Le choix de la durée d'observation, du fenêtrage et du traitement spectral doit être cohérent avec les phénomènes physiques recherchés. \textbf{En l'absence de fréquence d'échantillonnage confirmée, il serait incorrect d'attribuer des valeurs en hertz à des raies spectrales.} Toute exploitation spectrale physiquement interprétée sera donc conditionnée à la vérification de cette métadonnée.

\section{Principe des autoencodeurs}
Soit $\mathbf{x}\in\mathbb{R}^d$ un vecteur d'entrée. Un encodeur $f_{\theta}$ calcule une représentation latente $\mathbf{z}$ ; un décodeur $g_{\phi}$ produit sa reconstruction :
\begin{equation}
  \mathbf{z}=f_{\theta}(\mathbf{x}),\qquad
  \widehat{\mathbf{x}}=g_{\phi}(\mathbf{z}).
\end{equation}
Un objectif de référence est la minimisation de l'erreur quadratique moyenne :
\begin{equation}
  \mathcal{L}_{\mathrm{rec}}=\frac{1}{M}\sum_{j=1}^{M}
    \left\|\mathbf{x}^{(j)}-\widehat{\mathbf{x}}^{(j)}\right\|_2^2.
  \label{eq:aeloss}
\end{equation}
La forme de l'entrée (caractéristiques statistiques, fenêtres temporelles, représentations temps--fréquence) et l'architecture de l'autoencodeur devront être justifiées par les propriétés effectivement constatées des données.

\begin{figure}[htbp]
\centering
\begin{tikzpicture}[>=Latex, every node/.style={font=\sffamily\footnotesize},
    bloc/.style={draw=bleufonce!35,rounded corners=2mm,minimum height=13mm,
                 minimum width=24mm,align=center,fill=grisclair}]
  \node[bloc] (x) {Données\\d'entrée};
  \node[bloc,right=5mm of x,fill=bleupale] (enc) {Encodeur};
  \node[bloc,right=5mm of enc,fill=turquoise!13] (z) {Espace\\latent};
  \node[bloc,right=5mm of z,fill=bleupale] (dec) {Décodeur};
  \node[bloc,right=5mm of dec] (xh) {Données\\reconstruites};
  \draw[->,very thick,bleufonce] (x)--(enc);
  \draw[->,very thick,bleufonce] (enc)--(z);
  \draw[->,very thick,bleufonce] (z)--(dec);
  \draw[->,very thick,bleufonce] (dec)--(xh);
\end{tikzpicture}
\caption{Schéma conceptuel d'un autoencodeur (illustration réalisée pour ce rapport).}
\label{fig:autoencodeur}
\end{figure}

\section{Deux usages complémentaires}
\subsection{Regroupement non supervisé}
Les représentations latentes $\mathbf{z}$ peuvent être regroupées à l'aide de méthodes telles que $k$-means ou les mélanges gaussiens (GMM). La qualité des partitions sera évaluée indépendamment des étiquettes utilisées pour caractériser les acquisitions. Des approches par caractéristiques classiques et par réduction de dimension de type PCA pourront servir de comparateurs.

\subsection{Détection d'anomalies par reconstruction}
Un autoencodeur entraîné sur des données de fonctionnement sain peut fournir un score d'anomalie fondé sur l'erreur de reconstruction. L'hypothèse est qu'une observation suffisamment éloignée de la distribution saine sera moins bien reconstruite ; elle doit néanmoins être \textbf{vérifiée expérimentalement}, car un réseau trop flexible peut également reconstruire certains défauts. La variation des régimes de fonctionnement peut par ailleurs affecter les scores.

\section{Positionnement par rapport à la littérature}
Les travaux de Shao \emph{et al.} \cite{shao2017} explorent l'apprentissage de caractéristiques par autoencodeur pour le diagnostic de machines tournantes. Ceux de Malhotra \emph{et al.} \cite{malhotra2016} étudient une approche encodeur--décodeur appliquée à l'identification d'anomalies sur des séries temporelles multicapteurs. Ces études motivent l'exploration des autoencodeurs, sans permettre de présumer de leurs performances sur notre banc d'essai.

\acompleter{Construire une synthèse bibliographique critique : contexte de chaque article, entrées, architectures, protocole d'évaluation, limites et pertinence réelle pour le jeu de données. Ne retenir les techniques du manuel vibratoire qu'après examen de leur valeur ajoutée.}

\chapter{Présentation et qualification du jeu de données}
\label{chap:donnees}

\section{Origine et structure connue}
Les données proviennent d'un banc d'essai expérimental dédié à l'étude de comportements vibratoires de roulements. Elles sont rangées selon cinq états mécaniques et trois régimes désignés par les valeurs 10, 15 et 20. L'unité exacte de ces régimes doit encore être établie.

\begin{table}[htbp]
\centering\small
\caption{Catégories connues à ce stade.}
\label{tab:classes}
\begin{tabularx}{\textwidth}{@{}>{\ttfamily}lX@{}}
\toprule
\multicolumn{1}{@{}l}{\bfseries Dossier} & \bfseries Interprétation préliminaire\\
\midrule
SAFE & État considéré comme sain.\\
BALL & Défaut affectant un élément roulant.\\
INNER\_RACE & Défaut affectant la bague intérieure.\\
OUTER\_RACE & Défaut affectant la bague extérieure.\\
COMBINED & Défauts combinés ; composition exacte à confirmer.\\
\bottomrule
\end{tabularx}
\end{table}

Les fichiers sont au format \texttt{.xlsx} et contiennent notamment les deux colonnes \path{acceleration_vertbleu} et \path{acceleration_vert}. Ces noms désignent deux voies de mesure d'\textbf{accélération verticale} ; le terme « vert » ne doit pas être interprété comme une couleur de capteur. Les fichiers consultés ne comportent pas de colonne temporelle explicite.

\section{Informations métrologiques et expérimentales manquantes}
\begin{encadreattente}
\textbf{Points à confirmer avant de figer le protocole :}
\begin{itemize}
  \item fréquence d'échantillonnage et unités d'accélération ;
  \item signification physique des régimes 10, 15 et 20 ;
  \item position, synchronisation et calibration des deux capteurs ;
  \item durée des acquisitions, nombre d'échantillons et continuité éventuelle entre fichiers ;
  \item identité des essais indépendants et éventuelles répétitions ;
  \item nature, localisation et sévérité des défauts, notamment combinés ;
  \item correspondance entre les segments des noms de fichiers et les paramètres d'acquisition.
\end{itemize}
\end{encadreattente}

\section{Analyse exploratoire attendue}
L'analyse exploratoire devra produire au minimum un inventaire du nombre de fichiers et d'échantillons par état et par régime, la distribution des durées ou tailles d'acquisition, des contrôles de valeurs manquantes et de duplication, ainsi qu'une étude descriptive des amplitudes pour chacune des voies de mesure.

Les comparaisons graphiques devront distinguer les \textbf{effets de l'état mécanique} des \textbf{effets du régime}. Une éventuelle synchronisation des deux voies ne devra pas être présumée sans vérification.

\acompleter{Ajouter : schéma du banc d'essai, photographies et position des capteurs, tableau exhaustif de volumétrie, statistiques des données, traçabilité de la collecte, et résultats de l'analyse exploratoire (responsables techniques à préciser avec l'équipe).}

\chapter{Protocole expérimental envisagé}
\label{chap:methodes}
\begin{encadreinfo}
\textbf{Statut : protocole proposé, non encore validé.} Les choix finaux devront être justifiés à partir de l'analyse exploratoire, de la taille des données, des moyens de calcul et des consignes pédagogiques.
\end{encadreinfo}

\section{Construction des jeux d'entraînement, de validation et de test}
La séparation des données constitue une étape préalable à tout découpage en fenêtres et à tout ajustement de normalisation. Des fenêtres issues d'une même acquisition, voire de plusieurs fichiers contigus appartenant à un même essai, sont potentiellement corrélées. Les répartir au hasard entre entraînement et test pourrait conduire à une estimation artificiellement optimiste des performances.

Le principe retenu sera donc de \textbf{séparer des groupes d'acquisitions indépendants}, dès lors qu'une identification fiable de ces groupes est possible. Si l'identité des essais ne peut pas être établie, la limite devra être documentée et les conclusions adaptées. Les transformations apprises, notamment la normalisation et la sélection de variables, seront ajustées uniquement sur l'entraînement.

\section{Préparation des signaux}
Les opérations qui pourront être étudiées sont :
\begin{enumerate}
  \item contrôle d'intégrité des fichiers et traitement explicite des valeurs manquantes ;
  \item analyse des différences d'échelle entre les capteurs ;
  \item segmentation en fenêtres de taille justifiée, avec ou sans recouvrement ;
  \item standardisation calculée sur l'ensemble d'entraînement uniquement ;
  \item expérimentation d'entrées temporelles brutes, de caractéristiques vibratoires ou, si la fréquence d'échantillonnage est validée, de représentations spectrales.
\end{enumerate}
La taille des fenêtres sera d'abord spécifiée en \emph{nombre d'échantillons}. Sa conversion en durée physique n'est possible que lorsque $f_s$ est confirmé.

\section{Volet A : clustering}
Le premier volet visera à comparer une ou plusieurs représentations apprises par autoencodeur à des références telles que les caractéristiques statistiques classiques et/ou la réduction de dimension PCA. Les regroupements pourront être calculés par $k$-means et, si pertinent, par un modèle de mélange gaussien.

La qualité intrinsèque du regroupement pourra être examinée au moyen du coefficient de silhouette (lorsque ses conditions d'application sont remplies). Lorsque les étiquettes d'état sont disponibles et fiables, l'accord entre clusters et états connus pourra être étudié avec l'\emph{Adjusted Rand Index} (ARI) et la \emph{Normalized Mutual Information} (NMI). Les étiquettes ne devront pas servir à orienter artificiellement le choix final dans une expérience annoncée comme non supervisée.

\section{Volet B : détection d'anomalies}
Une approche consiste à entraîner un autoencodeur exclusivement sur des acquisitions \texttt{SAFE}, puis à calculer pour chaque fenêtre un score de reconstruction, par exemple
\begin{equation}
  s(\mathbf{x})=\frac{1}{d}\sum_{j=1}^{d}(x_j-\hat{x}_j)^2.
  \label{eq:score}
\end{equation}
Le seuil de décision devra être défini à partir d'un \textbf{jeu de validation séparé}, sans être ajusté sur le test final. Une méthode possible est un quantile fixé a priori de la distribution des scores de validation sains. Une alternative exploitant des défauts étiquetés pour la calibration devrait être explicitement qualifiée de protocole partiellement supervisé.

Les comparateurs pourront inclure un score de reconstruction PCA et Isolation Forest \cite{liu2008}. L'évaluation devra porter sur la précision, le rappel, le score F1, l'aire sous la courbe précision--rappel (PR-AUC) et les fausses alertes parmi les acquisitions saines, selon la pertinence de ces indicateurs pour la distribution des classes.

\section{Robustesse, biais et traçabilité}
Les résultats seront examinés \textbf{par régime de fonctionnement} et, si les métadonnées le permettent, \textbf{par voie de capteur}. Une séparation insuffisante entre les acquisitions ou une dépendance excessive au régime constitue un risque central.

Les expériences devront conserver les paramètres de prétraitement, l'identité des fichiers utilisés, les graines aléatoires, les hyperparamètres, les versions des bibliothèques et les métriques. En cas de multiples fenêtres provenant d'un même fichier, les indicateurs au niveau \emph{fichier ou essai} pourront compléter les métriques au niveau fenêtre, afin de ne pas traiter des observations corrélées comme indépendantes.

\section{Matrice expérimentale à confirmer}
\begin{table}[htbp]
\centering\small
\caption{Expériences prévues : liste indicative, à valider en équipe.}
\label{tab:experiences}
\begin{tabularx}{\textwidth}{@{}p{18mm}XXp{23mm}@{}}
\toprule
\bfseries ID & \bfseries Approche & \bfseries Objectif & \bfseries Statut\\
\midrule
C1 & Caractéristiques vibratoires + $k$-means & Référence clustering & À planifier\\
C2 & PCA + $k$-means & Référence réduction de dimension & À planifier\\
C3 & Autoencodeur + $k$-means / GMM & Clustering en espace latent & À planifier\\
A1 & Isolation Forest / PCA & Références de détection & À planifier\\
A2 & Autoencodeur entraîné sur SAFE & Score de reconstruction anormal & À planifier\\
\bottomrule
\end{tabularx}
\end{table}

\acompleter{Ajouter le protocole de découpage validé, la forme exacte des entrées, les modèles effectivement retenus, les espaces d'hyperparamètres, les métriques choisies et leur justification.}

\chapter{Expériences et résultats}
\label{chap:resultats}
\acompleter{Ce chapitre restera sans résultat chiffré tant que les expériences n'auront pas été exécutées et vérifiées. Ajouter successivement : environnement expérimental, vérification des découpages, courbes d'apprentissage, métriques de clustering, détection d'anomalies, robustesse par régime, comparaisons et analyses d'erreurs.}

\section{Caractérisation finale du jeu de données}
\acompleter{Présenter les effectifs par état, régime, capteur et jeu de données. Décrire toutes les exclusions ou corrections.}

\section{Comparaison des méthodes de clustering}
\acompleter{Insérer les métriques ARI/NMI/silhouette uniquement après calcul, avec les choix de $k$, la stabilité des résultats et les représentations visuelles appropriées.}

\section{Détection d'anomalies et analyse des faux positifs}
\acompleter{Reporter les courbes précision--rappel, matrices de confusion, seuils de validation et performances par régime. Décrire les anomalies mal détectées et les conditions des fausses alarmes.}

\section{Choix argumenté du ou des modèles}
\acompleter{Expliquer le compromis entre performance, robustesse, simplicité, coût de calcul et capacité d'intégration. Le modèle retenu n'est pas fixé à ce stade.}

\chapter{Architecture et démonstration de l'application}
\label{chap:application}
Ce chapitre documentera les composants développés pour rendre les expériences reproductibles et présenter le fonctionnement du système. Son contenu dépendra des choix effectivement validés par l'équipe.

\section{Architecture cible et circulation des données}
\acompleter{Décrire le dépôt, les notebooks, les artefacts modèles, la préparation des données, le backend et le frontend. Insérer un diagramme original de l'architecture seulement lorsqu'elle sera stabilisée.}

\section{Parcours utilisateur et limites de la démonstration}
\acompleter{Préciser les données acceptées, les contrôles, les résultats affichés et le sens des alertes. Ne pas présenter la démonstration comme un système de surveillance temps réel ou de maintenance industriellement validé sans essais correspondants.}

\section{Reproductibilité et déploiement}
\acompleter{Expliquer les dépendances, les instructions d'installation, l'export des modèles, les tests d'intégration et les limites de fonctionnement.}

\chapter{Discussion générale et conclusion}
\label{chap:discussion}

\section{Interprétation physique des résultats}
\acompleter{Mettre en relation les résultats des modèles avec les caractéristiques des signaux sans transformer une simple corrélation statistique en preuve d'un mécanisme mécanique. Revenir aux indicateurs vibratoires et aux effets de régime.}

\section{Limites scientifiques et opérationnelles}
\acompleter{Discuter : représentativité du banc, qualité des métadonnées, indépendance des essais, décalage de domaine, sensibilité des seuils, variations de vitesse, interprétabilité et absence éventuelle de suivi temporel de dégradation.}

\section{Conclusion et perspectives}
\acompleter{La conclusion sera rédigée après les expériences. Elle devra répondre à la problématique à partir des résultats mesurés, indiquer ce que l'étude ne permet pas de conclure et proposer des perspectives proportionnées aux preuves disponibles.}


\clearpage
\addcontentsline{toc}{chapter}{Références bibliographiques}
\bibliographystyle{unsrt}
\bibliography{references}

\appendix
\chapter{Suivi des versions et décisions}
\label{ann:suivi}
Ce journal soutient la documentation évolutive du projet et doit être actualisé lors des choix qui influencent les conclusions scientifiques.

\section{Historique des versions}
\begin{longtable}{@{}p{22mm}p{33mm}p{88mm}@{}}
\toprule
\bfseries Version & \bfseries Date & \bfseries Modification\\
\midrule
\endhead
0.1 & 27/10/2026 & Mise en place de la structure du rapport ; rédaction préliminaire du contexte, du cadre scientifique, de la description du jeu de données et du protocole proposé.\\
\bottomrule
\end{longtable}

\section{Registre des points à clarifier}
\begin{longtable}{@{}p{14mm}p{89mm}p{40mm}@{}}
\toprule
\bfseries ID & \bfseries Point à clarifier & \bfseries Statut\\
\midrule
\endhead
Q1 & Fréquence d'échantillonnage et unités des signaux & À confirmer\\
Q2 & Unités et signification physique des régimes 10/15/20 & À confirmer\\
Q3 & Structure des essais et indépendance des fichiers & À confirmer\\
Q4 & Positionnement et synchronisation des deux capteurs & À confirmer\\
Q5 & Définition exacte des défauts combinés & À confirmer\\
Q6 & Protocole final de segmentation et de validation & À décider après EDA\\
\bottomrule
\end{longtable}

\section{Registre des décisions scientifiques}
\acompleter{À chaque décision majeure, renseigner : date, question, options examinées, justification, données ou sources utilisées, auteur de la proposition, validation collective et conséquences pour les expériences.}

\end{document}
